\section*{第 \ref{ch:g12:selection-drift-speciation} 章——选择、漂变与物种形成}

\begin{solution}{exo:g12:selection-drift-speciation:1}
某\omterm{def:g10:universal-dna:gene}{基因}在\omterm{def:g12:selection-drift-speciation:population}{种群}全部拷贝中属于某一给定\omterm{def:g10:universal-dna:gene}{等位基因}的比例。
设 $p = 0.8$，$q = 0.2$，则 $AA$ 占 64\%，$Aa$ 占 32\%，$aa$ 占 4\%。
\end{solution}

\begin{solution}{exo:g12:selection-drift-speciation:2}
选择：不同\omterm{def:g10:universal-dna:gene}{等位基因}的携带者之间繁殖能力存在差异，
由环境决定，使频率朝某一方向改变。
漂变：有限\omterm{def:g12:selection-drift-speciation:population}{种群}抽样引起的频率随机涨落，
无方向性。选择取决于\omterm{def:g10:universal-dna:gene}{等位基因}的作用，
漂变则不然。
\end{solution}

\begin{solution}{exo:g12:selection-drift-speciation:3}
1880 年约为 75\%，当时煤烟已把树干熏黑；
1980 年约为 30\%，清洁空气法令实施后树干重新变浅。
\end{solution}

\begin{solution}{exo:g12:selection-drift-speciation:4}
下一代的\omterm{def:g10:universal-dna:gene}{等位基因}是上一代的一个样本；
样本越小，与来源的偏差越大，因此频率波动越剧烈，
某一\omterm{def:g10:universal-dna:gene}{等位基因}可能因偶然因素而丢失。
\end{solution}

\begin{solution}{exo:g12:selection-drift-speciation:5}
两个\omterm{def:g12:selection-drift-speciation:population}{种群}之间\omterm{def:g10:universal-dna:gene}{基因}交流中断，
且持续足够长的时间，使二者分化至不能再相互交配。
\end{solution}

\begin{solution}{exo:g12:selection-drift-speciation:6}
$a$ 的拷贝数：$160 + 40 = 200$（共 1000 个），故 $q = 0.2$，$p = 0.8$。
期望值：$0.64 \times 500 = 320$，$0.32 \times 500 = 160$，$0.04
\times 500 = 20$——与普查结果完全一致：\omterm{def:g12:selection-drift-speciation:population}{种群}处于
参考状态。
\end{solution}

\begin{solution}{exo:g12:selection-drift-speciation:7}
$q^2 = 0.0001$：万分之一的个体表现该性状；$2pq \approx
0.02$：五十分之一的个体携带该性状。99\% 的拷贝位于
杂合子中，选择对这些个体的该性状无从施压。
\end{solution}

\begin{solution}{exo:g12:selection-drift-speciation:8}
两条路径都在第 12 代达到终点：一条固定为 100\%，另一条
丢失。第三条路径在 20 代内游走，始终未触及任何一个边界
——同样是偶然因素所致；若时间更长，它终将触及。
\end{solution}

\begin{solution}{exo:g12:selection-drift-speciation:9}
浅色蛾在未受污染的乡村继续繁殖，在那里它们更占优势，
而蛾能飞行：迁入者每一代都把浅色\omterm{def:g10:universal-dna:gene}{等位基因}带回
煤烟区，因此该\omterm{def:g10:universal-dna:gene}{等位基因}从未消失。
\end{solution}

\begin{solution}{exo:g12:selection-drift-speciation:10}
20 个个体至多携带每个\omterm{def:g10:universal-dna:gene}{基因} 40 个拷贝，
只是原始\omterm{def:g10:universal-dna:gene}{等位基因}的极小样本；大部分在瓶颈效应中丢失，
随后小\omterm{def:g12:selection-drift-speciation:population}{种群}中的漂变又丢失了更多。\omterm{def:g12:selection-drift-speciation:population}{种群}数量得以恢复，
而\omterm{def:g10:universal-dna:gene}{等位基因}只能依靠新的\omterm{def:g11:mutations:mutation}{突变}重新出现。
\end{solution}

\begin{solution}{exo:g12:selection-drift-speciation:11}
十个\omterm{def:g12:selection-drift-speciation:population}{种群}在相同的环境变化之后朝同一方向变化，
并有合理的机制：选择。一个仅有 30 个个体的小\omterm{def:g12:selection-drift-speciation:population}{种群}，
没有报告任何原因：在机制得到揭示之前，
漂变是默认的解释。
\end{solution}

\begin{solution}{exo:g12:selection-drift-speciation:12}
$AA$：$0.85^2 \approx 72\%$，$Aa$：$2 \times 0.85 \times 0.15 \approx
26\%$，$aa$：$0.15^2 \approx 2\%$。$aa$ 个体死亡，使 $a$ 的拷贝减少；
但 $Aa$ 个体约占\omterm{def:g12:selection-drift-speciation:population}{种群}的四分之一，它们比 $AA$ 更能抵御疟疾，
并将其所携带的 $a$ 拷贝传给后代。$aa$ 的损失
被 $Aa$ 的优势所抵消，$a$ 保持在 15\%。
\end{solution}

\begin{solution}{exo:g12:selection-drift-speciation:13}
按颜色选择配偶是一种行为隔离。在浑浊的水中，
这一隔离失效，两个\omterm{def:g10:biodiversity-scales:species}{物种}发生杂交，因此尚未
积累遗传上的不相容性：这种分离近期才发生，
仍然是可逆的。
\end{solution}

\begin{solution}{exo:g12:selection-drift-speciation:14}
沿环状分布时，每个\omterm{def:g12:selection-drift-speciation:population}{种群}都与相邻\omterm{def:g12:selection-drift-speciation:population}{种群}相互交配，
相邻\omterm{def:g12:selection-drift-speciation:population}{种群}之间仅有微小差异；这些差异
沿环逐步累积，直到两端相遇时已不能相互交配。
整个渐变群同时可见：\omterm{prop:g12:selection-drift-speciation:speciation}{物种形成}是
小步累积的结果。两端之所以被称为\omterm{def:g10:biodiversity-scales:species}{物种}，
是因为相遇时它们表现得如同\omterm{def:g10:biodiversity-scales:species}{物种}：彼此之间没有\omterm{def:g10:universal-dna:gene}{基因}交流。
\end{solution}

\begin{solution}{exo:g12:selection-drift-speciation:15}
选择使某些\omterm{def:g10:universal-dna:gene}{等位基因}的携带者\emph{在此时
此地}繁殖得更多：与\omterm{def:g12:selection-drift-speciation:population}{种群}中其他成员相比，
它们更能在当前环境中留下后代，但这并不意味着它们在任何绝对意义上更好。
环境改变时，方向随之反转（如蛾的例子）；
漂变使\omterm{def:g12:selection-drift-speciation:population}{种群}出现差异，却与适应无关；
性选择则青睐妨碍生存的装饰性特征。
演化没有改进的方向，只有局部的筛选。
\end{solution}

\begin{solution}{pb:g12:selection-drift-speciation:1}
\textbf{1.} $D$ 的拷贝数：$2 \times 720 + 960 = 2400$；$d$ 的拷贝数：$960 +
2 \times 320 = 1600$；共 4000 个，故 $p = 0.6$，$q = 0.4$。

\textbf{2.} 期望值 $0.36 \times 2000 = 720$，$0.48 \times 2000 =
960$，$0.16 \times 2000 = 320$：与普查结果完全一致。是的。

\textbf{3.} $960/1680 \approx 57\%$。

\textbf{4.} 大多数 $d$ 拷贝位于深色携带者中。移除后：
1680 只老鼠，$D$ 拷贝 2400，$d$ 拷贝 960：$q = 960/3360 \approx
0.29$——较 0.40 下降，但远未消除。

\textbf{5.} 体型大：小岛\omterm{def:g12:selection-drift-speciation:population}{种群}规模小，发生漂变，故
频率因偶然而波动（且大陆迁入缺失，
\omterm{def:g11:mutations:mutation}{突变}可忽略）。

\textbf{6.} $DD$ 360，$Dd$ 480，$dd$ 320：共 1160 只繁殖个体。

\textbf{7.} $D$：$720 + 480 = 1200$；$d$：$480 + 640 = 1120$；共 2320 个：
$p \approx 0.52$，$q \approx 0.48$。

\textbf{8.} 下一代（每 2000 个）：$DD$ $0.52^2 \approx 27\%$
（535），$Dd$ $\approx 50\%$（999），$dd$ $\approx 23\%$（466）。经猫头鹰
捕食后为 267，500，466。$D$：$534 + 500 = 1034$；$d$：$500 + 932 =
1432$；$q \approx 0.58$。

\textbf{9.} 当 $q = 0.58$ 时：$DD$ 18\%（353），$Dd$ 49\%（974），$dd$
34\%（673）；存活个体 176，487，673；$d$：$487 + 1346 = 1833$，共 2672 个：
$q \approx 0.69$。趋势：0.40，0.48，0.58，0.69——每代约上升
0.1。

\textbf{10.} 每个 $D$ 携带者都是深色且暴露于猫头鹰之下，故
每代都对每个 $D$ 拷贝施加选择，$q$ 持续
稳步上升。\omterm{def:g11:genes-and-disease:genetic}{隐性}\omterm{def:g10:universal-dna:gene}{等位基因}一旦稀少，便隐藏于杂合子中，
而\omterm{def:g11:genes-and-disease:genetic}{显性}\omterm{def:g10:universal-dna:gene}{等位基因}从不隐藏：$D$ 可被驱至零。

\textbf{11.} $D$：$4 + 3 = 7$；$d$：$3 + 2 = 5$；$q = 5/12 \approx
0.42$，因抽样的偶然，接近小岛的 0.40。

\textbf{12.} 繁殖个体中 $D$ 有 5 个拷贝，$d$ 有 1 个：下一代的 $q$ 平均为 $1/6$，
但这唯一的 $d$ 拷贝可能无一传给幼体（$q = 0$），也可能传给多个；
$dd$ \omterm{def:g11:enzymes-and-phenotype:phenotype}{基因型}已经消失，
随之消失的还有大量变异。

\textbf{13.} 繁殖个体如此之少，
每一代哪些\omterm{def:g10:universal-dna:gene}{等位基因}得以传递都取决于偶然；
稀有\omterm{def:g10:universal-dna:gene}{等位基因}可能在一窝不幸的幼崽中就此消失。

\textbf{14.} $d$ 将发生漂变，并在若干代内被固定或
丢失——两种结局均出于偶然。一个邻近的小岛，
从相同起点出发，最终可能携带相反的\omterm{def:g10:universal-dna:gene}{等位基因}：两岛
在毛色上的差异并无适应性原因。

\textbf{15.} \omterm{prop:g12:selection-drift-speciation:drift}{遗传漂变}；群体规模极小。

\textbf{16.} 是的：它们很少与岛上的老鼠杂交，
且杂交后代不育——即生殖隔离。

\textbf{17.} 大海构成地理隔离；
在不同环境（无猫头鹰、其他条件）中的选择，
以及小\omterm{def:g12:selection-drift-speciation:population}{种群}中的\omterm{prop:g12:selection-drift-speciation:drift}{遗传漂变}。

\textbf{18.} 在不同季节繁殖的老鼠从不作为配偶相遇：
即使没有任何地理隔离，
\omterm{def:g10:universal-dna:gene}{基因}也不再流动，两个\omterm{def:g12:selection-drift-speciation:population}{种群}仍将持续分化。

\textbf{19.} 目前还不会：它们因行为和不育而相互隔离，
将作为两个并存的\omterm{def:g10:biodiversity-scales:species}{物种}存续。仅经过 100 年，它们
仍会自由杂交，并重新合并为一个\omterm{def:g12:selection-drift-speciation:population}{种群}。

\textbf{20.} 经过三代猫头鹰的选择，$q$ 从 0.40 升至约 0.69；
在小岛上，$d$ 因\omterm{prop:g12:selection-drift-speciation:drift}{遗传漂变}而被固定或丢失；
最终形成两个\omterm{def:g10:biodiversity-scales:species}{物种}。
\end{solution}
